\section{Main finite-game results}
\label{sec:finite-results}

The class partition and the equilibrium correspondence read different information from the same valuation map. Classes depend on payoff signs at a profile; pure best responses depend on unilateral payoff comparisons. This section states the two main results before introducing the benchmark parameterization.

For a contextual game $G_C=(C,A,B,V_C)$, write $V_C(a,b)=(X_C(a,b),Y_C(a,b))$. The best-response correspondences are
\[
\BR^A_C(b)=\argmax_{a\in A}X_C(a,b),
\qquad
\BR^B_C(a)=\argmax_{b\in B}Y_C(a,b).
\]
A profile $(a,b)$ is a pure Nash equilibrium if $a\in\BR^A_C(b)$ and $b\in\BR^B_C(a)$. It is strict if each action is the unique best response to the other player's action. We write $\NE(G_C)$ for the pure-equilibrium set.

\begin{definition}[Pareto and welfare dominance]
For two profiles $\sigma,\tau$ in a fixed context $C$, write $V_C(\sigma)=(x_\sigma,y_\sigma)$ and $V_C(\tau)=(x_\tau,y_\tau)$. Profile $\sigma$ Pareto-dominates $\tau$ if $x_\sigma\ge x_\tau$ and $y_\sigma\ge y_\tau$, with at least one strict inequality. It welfare-dominates $\tau$ if $W(V_C(\sigma))>W(V_C(\tau))$.
\end{definition}

\subsection{Exact preservation radius of the pure-equilibrium set}

For a profile $(i,j)\in A\times B$, collect all unilateral payoff advantages in
\[
\Delta_{ij}=
\{X(i,j)-X(i',j):i'\in A\setminus\{i\}\}
\cup
\{Y(i,j)-Y(i,j'):j'\in B\setminus\{j\}\}.
\]
Thus $(i,j)$ is a pure Nash equilibrium exactly when every element of $\Delta_{ij}$ is nonnegative.

\begin{theorem}[Exact pure-equilibrium-set robustness in finite games]
\label{thm:ne-radius}
Let $G$ be a finite two-player game in which each player has at least two actions. For each profile $(i,j)$ define
\[
r_{ij}=\begin{cases}
\tfrac12\min\Delta_{ij}, & (i,j)\in\NE(G),\\[2mm]
\tfrac12\max\{-d:d\in\Delta_{ij},\ d<0\}, & (i,j)\notin\NE(G).
\end{cases}
\]
Then
\[
\rho_{\NE}(G)=\min_{i,j}r_{ij}
\]
is the supremal open sup-norm radius that preserves the complete pure-equilibrium set: every payoff perturbation of norm less than $\rho_{\NE}(G)$ leaves $\NE(G)$ unchanged, while for every $\varepsilon>\rho_{\NE}(G)$ some perturbation of norm less than $\varepsilon$ changes it.
\end{theorem}

\begin{proof}
Each unilateral advantage is a difference of two payoff coordinates, so a sup-norm perturbation of radius $\delta$ changes it by less than $2\delta$. For an equilibrium profile, every advantage remains nonnegative when $\delta<\tfrac12\min\Delta_{ij}$. For a non-equilibrium profile, choose a negative advantage $d_*$ with largest absolute value. It remains negative when $\delta<-d_*/2$, so the profile remains outside the equilibrium set below $r_{ij}$. Taking the minimum over the finite profile set proves preservation.

For sharpness, choose a profile attaining the minimum. If it is an equilibrium, choose an advantage attaining $\min\Delta_{ij}$, lower its current payoff and raise the corresponding deviation payoff by a common amount strictly between $r_{ij}$ and $\varepsilon$; that advantage becomes negative. If it is not an equilibrium, increase the current payoff of each player by $r_{ij}$ and decrease that player's payoff at every profitable unilateral deviation by $r_{ij}$. Every initially negative advantage $d$ becomes $d+2r_{ij}\ge0$, so the profile becomes an equilibrium under a perturbation of norm $r_{ij}<\varepsilon$.
\end{proof}

The class-map radius developed in Section~\ref{sec:robustness} depends instead on distances to the payoff axes. The two radii can therefore differ even though both are computed from the same payoff vector.

\subsection{Configuration of intelligent and malicious equilibria}
\label{sec:general-theorem}

If $|A|=m$ and $|B|=n$, identify a game with its payoff vector $v\in\R^{2mn}$. A property is \emph{robust} when it holds on an open neighborhood of $v$.

\begin{lemma}[Indifference forced by coexisting equilibria]
\label{lem:coexist}
Let a finite two-player contextual benefit--loss game have two distinct pure Nash equilibria. If they share the actor's action, the affected-side player is indifferent between their two affected-side actions. If they share the affected side's action, the actor is indifferent between their two actor actions.
\end{lemma}

\begin{proof}
Suppose $(a,b)$ and $(a,b')$ are equilibria with $b\ne b'$. Optimality gives $Y_C(a,b)\ge Y_C(a,b')$ and $Y_C(a,b')\ge Y_C(a,b)$, hence equality. The shared-affected-action case is symmetric: optimality at $(a,b)$ and $(a',b)$ gives $X_C(a,b)=X_C(a',b)$.
\end{proof}

\begin{theorem}[Finite-game configuration theorem]
\label{thm:configuration}
Let $\sigma_I$ be an intelligent pure Nash equilibrium of a finite two-player contextual game and let $\sigma_M\ne\sigma_I$ be a malicious pure Nash equilibrium. Exactly one of the following relative positions occurs.
\begin{enumerate}[label=(\alph*)]
\item \textbf{Shared actor action: impossible.} The profiles cannot share the actor's action.
\item \textbf{Shared affected-side action: forced knife-edge.} If they share the affected side's action, then $X_C(\sigma_I)=X_C(\sigma_M)>0$, and $\sigma_I$ Pareto- and welfare-dominates $\sigma_M$. Games supporting any such pair lie in a finite union of proper hyperplanes and have empty interior.
\item \textbf{No shared action: the only possible robust regime.} Otherwise the profiles differ in both actions. This position alone is not sufficient for robustness. If both equilibria are strict and $W(\sigma_I)>W(\sigma_M)$, their class, equilibrium, and welfare conditions persist on an open neighborhood. Such instances exist already in $2\times2$ games.
\end{enumerate}
Thus sharing no action is necessary but not sufficient for robust coexistence of an intelligent equilibrium and a welfare-dominated malicious equilibrium. In a $2\times2$ game this is exactly the diagonal configuration.
\end{theorem}

\begin{proof}
If the profiles shared the actor's action, Lemma~\ref{lem:coexist} would imply $Y_C(\sigma_I)=Y_C(\sigma_M)$, contradicting $Y_C(\sigma_I)>0>Y_C(\sigma_M)$.

If they share the affected side's action, the lemma gives $X_C(\sigma_I)=X_C(\sigma_M)>0$. Together with $Y_C(\sigma_I)>0>Y_C(\sigma_M)$, this yields Pareto and welfare dominance. For a fixed profile pair the equality lies in a proper hyperplane of $\R^{2mn}$. There are finitely many profile pairs, so games admitting any such pair lie in a finite union of proper hyperplanes and have empty interior.

In the remaining case the profiles share neither action. Strict best-response, sign, and welfare-dominance conditions are finitely many strict linear inequalities in $v$; their intersection is open. Proposition~\ref{prop:diagonal-robust} supplies a nonempty instance.
\end{proof}

\begin{proposition}[Robust welfare-dominated malicious equilibria exist]
\label{prop:diagonal-robust}
Let $\Gamma^\diamond$ be the $2\times2$ contextual game
\[
V(G,U)=(1,4),\quad V(G,R)=(-2,2),\quad
V(P,U)=(-1,-3),\quad V(P,R)=(3,-1).
\]
Then $(G,U)$ is an intelligent strict Nash equilibrium, $(P,R)$ is a malicious strict Nash equilibrium, and $(G,U)$ welfare-dominates but does not Pareto-dominate $(P,R)$. All these properties hold for every payoff matrix at sup-norm distance strictly less than $3/4$ from $\Gamma^\diamond$. The bound is sharp for this certificate.
\end{proposition}

\begin{proof}
The strict best-response gaps at $(G,U)$ are $2$ and $2$, and those at $(P,R)$ are $5$ and $2$. The required class signs hold. Moreover,
\[
W(G,U)-W(P,R)=5-2=3>0,
\]
whereas $X(P,R)-X(G,U)=2>0$, so Pareto dominance fails. Each condition is a strict linear form $\ell(v)>0$. Under a sup-norm perturbation of size $\delta$, its value changes by at most $\delta\|\ell\|_1$. The smallest ratio $\ell(\Gamma^\diamond)/\|\ell\|_1$ is the welfare slack $3/4$. Hence all conditions persist for $\delta<3/4$. At $\delta=3/4$, decreasing both coordinates of $V(G,U)$ and increasing both coordinates of $V(P,R)$ by $3/4$ makes welfare equal, proving sharpness for the stated certificate.
\end{proof}
